\section{Experimental Setup}
\label{sec:setup}

The official test set will only be released in January 2027, so we use the public dev split as our test set and never use it for model selection. Instead, we hold out 10\% of the \emph{training} source pairs as a validation slice (the same pairs in every track, so there is no cross-lingual leakage). On this slice we select the learning rate from $\{2,5\}\cdot10^{-5}$ (on EN) and the best epoch (out of 10 for monolingual and 4 for mixed training). Selection always uses independent decoding, so it is blind to the decoder. We use AdamW, batch size 32 and 10\% warm-up, and we freeze the word embeddings to fit into 18\,GB of laptop memory (Appendix~\ref{app:hparams}). Every configuration is run with three seeds. All scores come from the official scorer.

\paragraph{Resources.} We use only open encoder checkpoints and the official training data.\footnote{\url{https://inigolopezgazpio.net/SemEval-2027-Task-2-DiCo-NLI/}} The only external data are the NLI corpora behind the mDeBERTa-NLI checkpoint (XNLI and translated English NLI). PhrasIS is not used.

\section{Results}
\label{sec:results}

\begin{table*}[t]
\centering\small
\setlength{\tabcolsep}{4pt}
\resizebox{\textwidth}{!}{%
\begin{tabular}{lccccccccccccccc}
\toprule
 & \multicolumn{3}{c}{EN} & \multicolumn{3}{c}{ES} & \multicolumn{3}{c}{EU} & \multicolumn{3}{c}{Mixed} & \multicolumn{3}{c}{Average} \\
\cmidrule(lr){2-4} \cmidrule(lr){5-7} \cmidrule(lr){8-10} \cmidrule(lr){11-13} \cmidrule(lr){14-16}
System & F1 & Soft & Hard & F1 & Soft & Hard & F1 & Soft & Hard & F1 & Soft & Hard & F1 & Soft & Hard \\
\midrule
Random & 25.3 & 18.1 & 5.5 & 25.3 & 18.1 & 5.5 & 25.3 & 18.1 & 5.5 & 25.1 & 18.4 & 6.2 & 25.2 & 18.1 & 5.7 \\
Constant \EQ & 11.0 & 100.0 & 31.4 & 11.0 & 100.0 & 31.4 & 11.0 & 100.0 & 31.4 & 11.0 & 100.0 & 31.4 & 11.0 & 100.0 & 31.4 \\
\midrule
XLM-R & 69.1 & 82.7 & 72.1 & 54.2 & 66.1 & 51.9 & 61.1 & 80.5 & 63.7 & 63.4 & 79.3 & 64.9 & 62.0 & 77.1 & 63.1 \\
\quad + joint decoding & 69.6 & 93.9 & 77.3 & 58.6 & \textbf{93.5} & 68.4 & 61.0 & \textbf{94.3} & 70.0 & 63.8 & 88.4 & 69.7 & 63.3 & \textbf{92.5} & 71.3 \\
\midrule
mDeBERTa & 79.6 & 87.1 & 81.0 & 77.3 & 85.6 & 78.3 & 76.5 & 82.7 & 75.7 & 73.4 & 81.8 & 73.8 & 76.7 & 84.3 & 77.2 \\
\quad + joint decoding & 80.3 & 91.8 & 84.5 & 77.9 & 89.3 & 81.1 & \textbf{76.7} & 87.5 & 78.8 & 73.5 & 86.2 & 76.6 & 77.1 & 88.7 & 80.2 \\
\midrule
mDeBERTa-NLI & 82.8 & 88.9 & 84.6 & \textbf{79.0} & 85.0 & 79.5 & 75.1 & 81.8 & 75.1 & 74.7 & 82.4 & 74.7 & 77.9 & 84.5 & 78.5 \\
\quad + joint decoding & 83.0 & 94.2 & 87.5 & 78.7 & 90.4 & 82.2 & 75.6 & 86.8 & 78.5 & 74.8 & 88.4 & 78.3 & 78.0 & 89.9 & 81.6 \\
\midrule
mDeBERTa + $\mathcal{L}_\text{cons}$ & 80.7 & 90.0 & 83.9 & 76.5 & 86.9 & 78.6 & 75.1 & 85.4 & 77.4 & 72.4 & 85.7 & 75.4 & 76.2 & 87.0 & 78.8 \\
\quad + joint decoding & 80.8 & 92.5 & 85.3 & 76.8 & 91.1 & 81.2 & 75.3 & 88.0 & \textbf{79.3} & 72.5 & 88.6 & 77.2 & 76.3 & 90.0 & 80.8 \\
\midrule
mDeBERTa-NLI 3-seed ens. & 83.7 & 89.9 & 86.6 & 78.8 & 87.4 & 80.5 & 76.4 & 82.7 & 77.3 & \textbf{76.0} & 83.9 & 76.6 & 78.7 & 86.0 & 80.2 \\
\quad + joint decoding & \textbf{84.3} & \textbf{94.9} & \textbf{89.9} & 78.9 & 91.3 & \textbf{82.7} & 76.2 & 85.9 & 79.1 & 75.9 & \textbf{88.9} & \textbf{79.5} & \textbf{78.9} & 90.3 & \textbf{82.8} \\
\midrule
\textit{mDeBERTa + struct (O2)} & \textit{91.2} & \textit{100.0} & \textit{89.5} & \textit{90.2} & \textit{100.0} & \textit{88.2} & \textit{88.7} & \textit{100.0} & \textit{86.5} & \textit{87.9} & \textit{100.0} & \textit{85.6} & \textit{89.5} & \textit{100.0} & \textit{87.5} \\
\bottomrule
\end{tabular}}
\caption{Dev results (\%). Learned systems: mean over 3 seeds (std.\ in Appendix~\ref{app:std}); the \emph{3-seed ens.} rows are a single ensemble (averaged probabilities), not a mean. \emph{joint}: joint bidirectional decoding (Eq.~\ref{eq:joint}). Bold: best learned (non-diagnostic) value. The last row exploits observation O2 and is shown as a diagnostic only.}
\label{tab:main}
\end{table*}


\paragraph{Encoders.} Under independent decoding, mDeBERTa outperforms XLM-R by 14.7 F1 points on average (Table~\ref{tab:main}). XLM-R is also unstable on these small training sets: one Spanish seed collapsed to 41.3 validation F1. On English, mDeBERTa (79.6 F1 / 81.0 HardCons) is above the DeBERTa-v3 pilot baseline \citep[75 / 79;][]{apaolaza2026assessing}, although the two setups are not identical. Scores decrease from EN to ES to EU to Mixed. Within the mixed track, the Spanish--Basque combination is the hardest (mDeBERTa: 71.6 F1 vs.\ 75.1 for en-es).

\paragraph{Joint decoding.} Over all 72 evaluations (the four systems of Table~\ref{tab:main} and the two transfer settings of Table~\ref{tab:transfer}, $\times$ 3 seeds $\times$ 4 tracks), joint decoding improved HardCons \emph{every time}, by $+3.9$ points on average. The gain was $+3.0$ for mDeBERTa and $+8.2$ for XLM-R: the weaker and less stable the model, the larger the gain. It also raised SoftCons in every evaluation ($+6.5$) and F1 slightly ($+0.5$, positive in 51 of 72). Joint decoding is complementary to seed ensembling. Our best in-track system, a 3-seed ensemble of NLI-pretrained mDeBERTa with joint decoding, reaches 78.9 F1 and 82.8 HardCons on average. XLM-R's top SoftCons only reflects its rare \NO predictions (29\% recall; O1).

\paragraph{NLI pretraining.} Under independent decoding, starting from the NLI-tuned checkpoint helps on English ($+3.2$ F1) and Spanish ($+1.7$), and slightly on Mixed ($+1.3$). It does not help on Basque ($-1.4$), a language absent from its NLI training data. On average it adds $+1.4$ HardCons under joint decoding.

\paragraph{Training-time vs.\ decoding-time consistency.} The consistency loss lowers the share of inconsistent pairs (2.5$\rightarrow$1.7\%), raising HardCons by $+1.6$ and SoftCons by $+2.7$ under independent decoding, at a cost of $-0.5$ F1. However, joint decoding alone gives about twice that HardCons gain ($+3.0$), and adding the loss underneath joint decoding brings only $+0.6$ HardCons (80.2$\rightarrow$80.8) while still costing 0.8 F1. Enforcing consistency at decoding time is therefore both cheaper and more effective than learning it.

\begin{table}[t]
\centering\small
\setlength{\tabcolsep}{2.5pt}
\begin{tabular}{@{}llccccc@{}}
\toprule
Model & Dec. & OK & \NO & Inc. & Wrong & F1\textsubscript{\NO} \\
\midrule
XLM-R & indep & 63.1 & 9.9 & 13.0 & 14.0 & 33.1 \\
 & joint & 71.3 & 7.5 & 0.0 & 21.2 & 33.0 \\
mDeBERTa & indep & 77.2 & 13.2 & 2.5 & 7.1 & 53.9 \\
 & joint & 80.2 & 11.3 & 0.0 & 8.5 & 54.0 \\
\bottomrule
\end{tabular}
\caption{Outcome of gold reversible pairs (\% of pairs, dev, mean over tracks and seeds): both correct (OK), a direction predicted \NO, inconsistent (Inc.), or consistent but wrong. F1\textsubscript{\NO}: instance-level F1 of \NO.}
\label{tab:errors}
\end{table}


\paragraph{Where do errors come from?} Table~\ref{tab:errors} decomposes the outcome of every gold reversible pair. For mDeBERTa, SoftCons fails on 15.7\% of pairs, close to the 16\% of the pilot study, yet only 2.5\% are directionally inconsistent: the dominant failure is predicting \NO for a reversible pair (13.2\%). The \NO class reaches only 54 F1. Joint decoding removes all inconsistencies and also repairs some spurious \NO predictions (13.2$\rightarrow$11.3\%), because a confident direction overrides the other one. The diagnostic \emph{struct} row in Table~\ref{tab:main} quantifies how much the \NO decision matters. Using the pair structure (O2) instead of the model to detect \NO raises average F1 from 77.1 to 89.5 and HardCons from 80.2 to 87.5 over joint decoding. On public data, the hardest part of the task is thus deciding \emph{whether} a relation is reversible, not \emph{in which direction}, and, unless the test data hides the structure, also the easiest to shortcut.

\begin{table}[t]
\centering\small
\setlength{\tabcolsep}{4pt}
\begin{tabular}{@{}lcccc@{}}
\toprule
Training data & EN & ES & EU & Mixed \\
\midrule
In-track & 84.5 & 81.1 & 78.8 & 76.6 \\
EN only & 84.0 & 80.7 & 68.5 & 54.4 \\
Pooled & 84.8 & 83.6 & 81.5 & 77.7 \\
\bottomrule
\end{tabular}
\caption{Dev HardCons (\%) of mDeBERTa with joint decoding (3 seeds) by training data: the track's own (in-track), Track~1 only (zero-shot on the others), or all four tracks (pooled).}
\label{tab:transfer}
\end{table}


\paragraph{Cross-lingual transfer.} A model trained on English only keeps its HardCons on Spanish (80.7 vs.\ 81.1 in-track; Table~\ref{tab:transfer}) but loses 3.5 F1 points. It predicts \NO less often (recall 42 vs.\ 59\%), which costs F1 on the \NO items but spares the reversible pairs. On Basque and on the mixed track it loses 10 and 22 HardCons points, so cross-lingual pairs require cross-lingual supervision. Pooling all four training sets gives the best single-model HardCons (81.9 vs.\ 80.2 in-track and 81.6 for mDeBERTa-NLI), with the largest gains on Basque ($+2.7$) and Spanish ($+2.5$). Its 3-seed ensemble (78.7 F1 / 82.9 HardCons) matches our best in-track system without external NLI data.

\section{Work Plan and Task Division}
\label{sec:plan}

All components of Section~\ref{sec:method} are implemented. For the final report (28 Dec) we plan to:
(i) combine NLI pretraining and pooled training in a 3-seed ensemble, and score the swapped input of \emph{every} instance with the model itself;
(ii) calibrate the \NO logit on the validation slice and add reversed \NO training items (also \NO) with loss weight 0.5;
(iii) tune $\lambda\in\{0.1, 0.3, 1\}$;
(iv) qualitative error analysis; optionally, a small open LLM (zero-shot, independent vs.\ joint decoding).

\begin{table}[h]
\centering\small
\setlength{\tabcolsep}{3pt}
\begin{tabular}{@{}lp{0.74\columnwidth}@{}}
\toprule
Member & Code / analysis; plan item; sections written \\
\midrule
K\"{u}\c{c}\"{u}k & Scorer integration, baselines, error analysis (Tab.~\ref{tab:errors}); (iv); \S1, \S6 \\
\c{S}ahin & Training pipeline, joint decoding, experiment runs; (i); \S4, \S5, \S7, \S8 \\
Okudurlar & Encoders, NLI transfer, consistency loss; (iii); \S3, \S6 \\
Tumurkhuu & Data artifacts (O1/O2), cross-lingual transfer (Tab.~\ref{tab:transfer}); (ii); \S2, \S6 \\
\bottomrule
\end{tabular}
\caption{Task division.}
\label{tab:plan}
\end{table}
